% Źródła: materialy/2026-10-05/notatki.pdf oraz tablica-01, -02, -03.jpeg.
\chapter{Aksjomatyczna definicja prawdopodobieństwa}
\label{chap:aksjomaty}
\poczatekwykladu{1}

\section{Przestrzeń probabilistyczna}
\label{sec:przestrzen}

\begin{definicja}
\emph{Przestrzenią probabilistyczną} nazywamy trójkę
$(\Omega,\F,\PP)$, gdzie:
\begin{enumerate}
  \item $\Omega$ jest niepustym zbiorem, zwanym \emph{przestrzenią
    stanów elementarnych};
  \item $\F$ jest rodziną podzbiorów $\Omega$ taką, że:
    \begin{align*}
      &\varnothing\in\F,\\
      &A\in\F\ \Longrightarrow\ A^c\in\F,\\
      &A_1,A_2,\ldots\in\F\ \Longrightarrow\
        \bigcup_{n=1}^{\infty}A_n\in\F;
    \end{align*}
  \item $\PP\colon\F\to[0,1]$ jest funkcją spełniającą warunki:
    \begin{align}
      \PP(\Omega)&=1,\label{eq:normalizacja}\\
      \PP\!\left(\bigcup_{n=1}^{\infty}A_n\right)
        &=\sum_{n=1}^{\infty}\PP(A_n),
        \label{eq:addytywnosc}
    \end{align}
    przy czym wzór \eqref{eq:addytywnosc} obowiązuje dla każdej
    rodziny $A_1,A_2,\ldots\in\F$ zdarzeń parami rozłącznych, czyli
    $A_i\cap A_j=\varnothing$ dla $i\ne j$.
\end{enumerate}
Rodzinę $\F$ nazywamy \emph{$\sigma$-ciałem}, a jej elementy
\emph{zdarzeniami}. Funkcję $\PP$ nazywamy \emph{prawdopodobieństwem}.
Warunek \eqref{eq:addytywnosc} to \emph{przeliczalna addytywność}.
\end{definicja}

\begin{uwaga}
Przyjmujemy następującą terminologię:
\begin{enumerate}
  \item $\omega\in\Omega$ jest stanem elementarnym.
  \item Zdarzenia $A,B\in\F$ \emph{wykluczają się}, jeżeli
    $A\cap B=\varnothing$.
  \item $A^c=\Omega\setminus A$ jest \emph{zdarzeniem przeciwnym}
    do zdarzenia $A$.
  \item $\varnothing$ jest \emph{zdarzeniem niemożliwym},
    a $\Omega$ jest \emph{zdarzeniem pewnym}.
\end{enumerate}
\end{uwaga}

\begin{fakt}
\label{fakt:pusty}
Prawdopodobieństwo zdarzenia niemożliwego wynosi zero:
\[
  \PP(\varnothing)=0.
\]
\end{fakt}

\section{Zbiory borelowskie i miara Lebesgue'a}
\label{sec:borel}

\begin{definicja}
\emph{$\sigma$-ciałem zbiorów borelowskich} na $\R$
(odpowiednio na $\R^2$ lub $\R^n$) nazywamy najmniejsze
$\sigma$-ciało zawierające wszystkie zbiory otwarte.
Oznaczamy je przez $\B(\R)$ (odpowiednio $\B(\R^2)$ lub $\B(\R^n)$).
Na prostej wystarczy rozważać otwarte przedziały $(a,b)$.
Elementy tego $\sigma$-ciała nazywamy \emph{zbiorami borelowskimi}.
\end{definicja}

Na przestrzeni $(\R,\B(\R))$ istnieje miara $\lambda$ taka, że
\[
  \lambda\bigl((a,b)\bigr)=b-a
  \qquad\text{dla }a<b.
\]
Jest to \emph{miara Lebesgue'a}. Mierzy ona długość i spełnia
warunek przeliczalnej addytywności:
\[
  \lambda\!\left(\bigcup_{n=1}^{\infty}A_n\right)
    =\sum_{n=1}^{\infty}\lambda(A_n)
\]
dla zbiorów borelowskich $A_n$ parami rozłącznych.

\begin{fakt}
\label{fakt:miara-zero}
Każdy zbiór jednopunktowy ma miarę Lebesgue'a równą zero:
\[
  \lambda(\{x\})=0\qquad\text{dla każdego }x\in\R.
\]
W szczególności $\lambda(\Q)=0$.
\end{fakt}

Istotnie, zbiór liczb wymiernych jest przeliczalny. Jeśli
$\Q=\{q_1,q_2,\ldots\}$, przy czym liczby $q_i$ są różne, to
\[
  \lambda(\Q)
    =\lambda\!\left(\bigcup_{i=1}^{\infty}\{q_i\}\right)
    =\sum_{i=1}^{\infty}\lambda(\{q_i\})
    =0.
\]

\section{Prawdopodobieństwo geometryczne}
\label{sec:geometryczne}

Niech $\Omega\subseteq\R$ będzie zbiorem borelowskim, dla którego
$0<\lambda(\Omega)<\infty$. Przyjmujemy
\[
  \F=\{A\cap\Omega:A\in\B(\R)\}.
\]
Dla każdego $B\in\F$ definiujemy
\begin{equation}
  \PP(B)=\frac{\lambda(B)}{\lambda(\Omega)}.
  \label{eq:geometryczne}
\end{equation}
W ten sposób otrzymujemy przestrzeń probabilistyczną
$(\Omega,\F,\PP)$.

Analogicznie, dla borelowskiego zbioru $\Omega\subseteq\R^2$
o dodatnim i skończonym polu, prawdopodobieństwo zdarzenia
$A\in\{B\cap\Omega:B\in\B(\R^2)\}$ określamy wzorem
\begin{equation}
  \PP(A)=\frac{\pole(A)}{\pole(\Omega)}.
  \label{eq:pole}
\end{equation}

\section{Przykład: suma dwóch losowanych liczb}
\label{sec:suma}

\begin{przyklad}
Losujemy niezależnie i jednostajnie dwie liczby z przedziału $[0,1]$.
Jakie jest prawdopodobieństwo, że ich suma jest większa od $1/2$?
\end{przyklad}

\noindent\emph{Rozwiązanie.}
Przestrzenią stanów elementarnych jest kwadrat jednostkowy
\[
  \Omega=[0,1]^2=\{(x,y):x,y\in[0,1]\}.
\]
Rozważamy zdarzenie
\[
  A=\{(x,y)\in[0,1]^2:x+y>\tfrac12\}.
\]
Jego dopełnienie $A^c$ jest trójkątem o wierzchołkach
$(0,0)$, $(1/2,0)$ i $(0,1/2)$, przedstawionym na
rysunku~\ref{fig:suma}. Stąd
\[
  \pole(A^c)=\frac12\cdot\frac12\cdot\frac12=\frac18,
  \qquad \pole(\Omega)=1.
\]
Korzystając ze wzoru \eqref{eq:pole}, otrzymujemy
\begin{equation}
  \PP(A)=1-\PP(A^c)
    =1-\frac{\pole(A^c)}{\pole(\Omega)}
    =1-\frac18=\boxed{\frac78}.
  \label{eq:wynik-suma}
\end{equation}

\begin{figure}[htbp]
  \centering
  \includegraphics[width=7.6cm]{suma-dwoch-liczb.pdf}
  \caption{Kwadrat $\Omega=[0,1]^2$ oraz zdarzenia
    $A=\{x+y>1/2\}$ i $A^c=\{x+y\leq1/2\}$.}
  \label{fig:suma}
\end{figure}

\begin{uwaga}
Na tablicy zapisano $x+y\geq1/2$. Odcinek $x+y=1/2$ ma pole
równe zero, więc zastąpienie znaku $>$ znakiem $\geq$
nie zmienia wyniku.
\end{uwaga}
